%! Solving Rational Equations — Unit 5, Lesson 5.6 — Exit Ticket (Answer Key)
\documentclass[10pt]{article}
\usepackage{algebra2-article}
\usepackage{algebra2-key}   % pulls in -boxes; do NOT also load -boxes
\newcommand{\TallMath}[1]{$\displaystyle #1\rule[-1.4em]{0pt}{3.2em}$}

\begin{document}

\pageheader{Unit 5, Lesson 5.6}{Exit Ticket --- Answer Key}
\namedateperiod

\vspace{0.06in}

No notes. Write your restriction list \emph{before} you touch the equation.

\vspace{0.04in}

\begin{enumerate}[leftmargin=1.8em, itemsep=9pt]
  \item \textbf{Solve completely.} \; $\dfrac{x}{x-1}+\dfrac{1}{x}=\dfrac{1}{x^2-x}$

        \vspace{3pt}
        \textbf{Step 1.} Factor: $x^2-x={\color{keyred}\mathbf{x(x-1)}}$. \;
        Restrictions: $x\neq$ \ans{$0,\;1$}. \; LCD $=$ \ans{$x(x-1)$}

        \vspace{3pt}
        \textbf{Step 2.} Multiply every term and cancel:
        \; \ans{$x\cdot x$} $+$ \ans{$(x-1)$} $=$ \ans{$1$}

        \vspace{3pt}
        \textbf{Step 3.} Solve: $0=x^2+{\color{keyred}\mathbf{1}}\,x-{\color{keyred}\mathbf{2}}
        =(x+{\color{keyred}\mathbf{2}})(x-{\color{keyred}\mathbf{1}})$. \; Candidates \ans{$-2,\;1$}

        \vspace{3pt}
        \textbf{Step 4.} Extraneous \ans{$1$} \quad Solution set \ans{$\{-2\}$}

  \item \textbf{Justify it.} In one or two sentences, explain to a classmate why the candidate you threw
        away in item 1 is not a solution --- and why that does \emph{not} mean anyone made an algebra
        mistake. Use the words \textbf{denominator} and \textbf{extraneous}.
        \ansline{$x=1$ makes the denominators $x-1$ and $x^2-x$ equal to zero, so the original
        equation is undefined there and $1$ cannot be a solution --- it is extraneous.}
        \ansline{No mistake was made: multiplying both sides by the LCD multiplies by zero at $x=1$,
        and that step cannot be undone, so the cleared equation has a root the original never had.}
        \ansline{Full credit needs both halves --- the zero denominator \emph{and} the reason the extra
        root appeared. ``It doesn't work when you plug it in'' alone is not a justification.}

  \item \textbf{SOL-style multiple choice.} What is the solution set of
        \[ \frac{x}{x+3}+\frac{3}{x-3}=\frac{18}{x^2-9}\;? \]
        \vspace{-8pt}
        \begin{enumerate}[label=\Alph*., leftmargin=1.9em, itemsep=1pt]
          \item $\{3\}$
          \item $\{-3\}$
          \item $\{3,\,-3\}$
          \item no solution \quad \textcolor{keyred}{\textbf{$\leftarrow$ correct}}
        \end{enumerate}
\end{enumerate}

\begin{teachernote}
\textbf{Item 1.} Clearing gives $x^2+(x-1)=1$, so $x^2+x-2=0$ and $(x+2)(x-1)=0$. The restriction list
is $\{0,1\}$, so $x=1$ is extraneous and the solution set is $\{-2\}$. Two things to watch: students who
factored $x^2-x$ as $x(x-1)$ but then wrote only $x\neq1$ have lost half the list (and would still get
this item right, which is why item 3 exists); and students who reject $-2$ for being negative have
imported a rule from the modeling problems that does not apply here.

\vspace{3pt}
\textbf{Item 2 is the A2.EI.4d item} and should be graded on the \emph{reason}, not the verdict.
Accept anything that names a zero denominator in the \emph{original} equation and connects the extra
root to the LCD step. Do not accept ``it doesn't check'' or ``it's extraneous'' with no mechanism ---
that restates the verdict.

\vspace{3pt}
\textbf{Item 3.} Clearing by $(x+3)(x-3)$ gives $x(x-3)+3(x+3)=18$, so $x^2=9$ and the candidates are
$3$ and $-3$ --- \emph{both} excluded. The answer is \textbf{D}. Each distractor is a specific failure:
\textbf{A} and \textbf{B} keep one candidate (usually whichever one the student happened to check),
and \textbf{C} skips Step 4 entirely. A student who picks C did the algebra perfectly.

\vspace{3pt}
\textbf{Sorting the pile.} ``Got it'' / ``no restriction list'' / ``list written but never used'' /
``rejected for the wrong reason.'' The middle two are the same fix and it is a fast one --- require the
list at the top of the page before any work. The last pile is the one to watch in 5.7, where rejecting
a solution \emph{for context} becomes correct and the two reasons have to stay separate.
\end{teachernote}

\end{document}
